\section{What case law can contribute to an interpretation}
\label{sec:legal-interpretation-reuse}

A product provider can retain responsibility for an outsourced arrangement
while supervising it poorly. The arrangement can also operate successfully
during a period for which the provider's authority is unclear. These
possibilities matter because a program that asks only whether
\texttt{responsibility\_retained} is true cannot discover which of these
questions its input was intended to answer. It can apply the supplied judgment
consistently while leaving the important legal distinction unexamined.

The tokenisation circular gives this problem a precise setting.
Paragraph~10 states that Product Providers should remain ultimately
responsible for the management and operational soundness of the arrangement
and ownership record keeping, regardless of outsourcing. Paragraph~14 says
they should “confirm and, upon SFC's request, demonstrate” the specified
matters. Footnote~5 says they should “confirm and demonstrate” compliance
with the preceding requirements and the specified assurance concerning smart
contracts. The scope and interaction of these statements must be interpreted;
neither a responsibility flag nor a general request flag settles them
\citep[paras.~10--16 and n.~5]{sfctoken2026}.

The current investigation consequently has two substantive tasks. It must
justify the proposed meaning of each duty, and it must explain what evidence
would establish the facts needed to apply that meaning. Case law may help
with both, by explaining legally significant distinctions and the reasons
they matter. To use a judgment responsibly, however, the product must identify
what the court decided, the facts on which that reasoning depended and the
conditions under which the reasoning bears on the present question. A
retrieved paragraph is the beginning of that account.

\subsection{The distinctions still missing from the product}

The product assessment of 5~October~2026 records 35 reading identities in two
roles. Those 70 role records are not 70 independently established meanings.
All 46 original challenge obligations remain open. The recorded 20
mechanical guards and 1,040 synthetic checks establish behavior under declared
premises; the selected repair package records no accepted substantive repair.
These counts describe that frozen assessment, rather than a fresh experiment
or a current legal error rate.\footnote{The assessment and exact package
identities are retained in
\srcpath{docs/implementation/assurance-evidence-master/product-gap-audit-2026-10-05.md}.}

There are also narrower defects that can be corrected without deciding the
law. Some operative explanations still say that referenced texts are
unavailable even though acquisition records show that they have been
obtained. Acquisition and applicability are different questions. Other fields
state the request condition categorically while a qualification elsewhere
preserves the unresolved interaction with footnote~5. The uncertainty must be
expressed where the duty is described, so a reader does not receive a settled
answer in one place and its withdrawal in another.

The more difficult gaps concern actors, activities, effective periods,
ownership, incidents and legal consequences. A provider's supervision is not
an intermediary's conduct; possession of a token does not by itself settle
legal title; and an incident followed by remediation is not the same factual
history as no incident. The word “should” also requires attention to the
authority and context in which it appears. Satisfying one encoded predicate
cannot establish overall compliance or the absence of civil liability.

These are reasons to improve the representation of meaning. They are not
reasons to require a unique answer where the available sources leave the
question open. A useful interpretation may identify two supported proposals,
show exactly where their consequences differ and state the missing premise
that prevents choosing between them.

\subsection{A precedent supplies a reason for a distinction}

HYPO provides an early, precise account of reasoning from cases. A proponent
identifies a precedent and the similarities that support its application. An
opponent distinguishes that precedent or supplies a more relevant contrary
case. The response must address the distinction rather than repeat the
original outcome. CATO develops hierarchies of legally relevant factors and
ways to emphasize or downplay differences. The benefit is an inspectable
argument about why a fact matters, rather than similarity between two blocks
of prose \citep[secs.~2--3]{hypolegacy2017}.

Suppose, solely as a design example, that one proposed interpretation of
responsibility turns on retained decision-making authority, while another
also requires evidence that supervision was exercised. A useful case
comparison would identify the authority and supervision findings in the
earlier case, the proposition for which that case is cited, and the reason a
difference in supervision matters here. It would not simply report that both
arrangements involved outsourcing or that both disputes had the same outcome.
The example supplies no real precedent and settles neither interpretation;
it identifies the information an actual precedent analysis must recover.

AGATHA extends this approach by constructing competing case-law theories.
Its inputs include cases represented through factors and outcomes, with
associated rules and value preferences. Search compares candidate theories
using such considerations as explanatory power, simplicity and completeness.
The formal representation makes rival explanations available for examination,
but the heuristic search need not find the best theory. The advocacy setting
also permits preference moves that cannot be imported as neutral findings
about governing law \citep[secs.~2--3, 6--8 and 10]{agatha2005}.

The limitation is consequential. These systems start after much of the
difficult interpretive work has been encoded. A factor such as retained
authority is not the raw text of a judgment. Someone or something has selected
the relevant passage, distinguished an allegation from a finding and decided
what proposition it supports. We can borrow the analogy, distinction and
counterexample operators while keeping that earlier justification visible.
The inspected material supports design reuse; it does not establish that an
official, reusable implementation of the legacy systems is available.

\subsection{Retrieving the right case is an incomplete success}

A case-law product must first find relevant material, and this is a useful
task in its own right. CLERC develops case retrieval and analysis-generation
tasks from a large US case-law collection. Its retrieval labels are derived
from citations in the original documents. Its generation study evaluates
text similarity and citation behavior, while its discussion explains why a
correct citation can accompany a deficient analytical proposition. A system
may cite the expected case but fail to state the proposition for which that
case is relevant \citep[secs.~3--5, especially sec.~5.2]{clerc2024}.

That distinction rules out citation accuracy as a substitute for an account
of the holding. CLERC also concatenates majority, concurring and dissenting
opinions in its case-document construction. That choice serves its retrieval
task, but a precedent interpreter must preserve the speakers and opinion
types. A dissent can explain a serious objection without supplying the
court's governing rationale. Procedural posture and later treatment can
further limit what an otherwise accurate quotation establishes
\citep[sec.~3.1]{clerc2024}.

The open-source \href{https://github.com/freelawproject/eyecite}{eyecite}
library can identify citations and connect shortened references to earlier
citation forms. Resolving them to actual judgments still requires a corpus
or lookup mechanism. Its principally US-oriented conventions need a
Hong Kong coverage check. It is not a service establishing that a case
remains good law. Similarly, a search that finds no contrary authority
establishes a search outcome within its coverage, not the absence of contrary
law.

Commercial tools do not remove the need for this distinction.
\citet[secs.~5--6]{legalrag2024} evaluated 202 queries against the versions of
three commercial research tools available during their 2024 study and
reported hallucination rates between 17\% and 33\% under their correctness
and grounding rubric. These are historical sample results, based on the
authors' human evaluation, not measurements of current 2026 products or
quality labels for LegalMath. Their relevance here is narrower: supplying
retrieved legal sources did not itself establish correct analysis.

\subsection{An argument can fail at its premise, rule or conclusion}

Once a proposed reading and its support are explicit, objections need more
structure than a single agree-or-disagree label. A challenge might say that
the source does not establish a factual premise. It might accept the premise
but contest the inference drawn from it. Or it might offer a contrary
conclusion supported by another rule. These differences determine what a
repair must change.

ASPIC+ distinguishes these attacks through structured arguments built from
premises and strict or defeasible rules. Undermining challenges a premise,
rebutting challenges a conclusion, and undercutting challenges the use of an
inference rule. Preferences can determine whether an attack succeeds as a
defeat. This provides a useful vocabulary for the paragraph-14 dispute:
one argument may challenge the proposed scope of the request condition,
while another challenges the priority or applicability assigned to the
footnote \citep[secs.~2--3]{aspic2010}.

The semantics then asks which arguments can jointly survive the declared
attacks. Grounded semantics yields one extension; preferred semantics can
yield several. Here an extension is a set of arguments accepted under the
specified formal definitions. It is not a set of independently established
legal truths. The premises, contrary relations, inference rules and
preferences remain part of what must be justified.

The inspected \href{https://github.com/DaphneOdekerken/PyArg}{PyArg}
implementation is a practical Python candidate for a bounded prototype.
Its conversion to an abstract argumentation framework chooses a last-link
elitist ordering when no ordering is supplied. That default must be made
explicit. A software preference cannot decide which reading the SFC intended.
Argument generation must also be bounded, with any unfinished frontier
reported. A truncated set of arguments cannot establish that no objection
exists.

ASPIC+'s rationality results require their particular assumptions. The
author's correction to the well-formedness definition is therefore a material
dependency for any attempted proof reuse. Our proposed use begins with
explicit argument and attack definitions; it does not import every theorem
for an arbitrary generated graph. A separate independent check must compare
the chosen implementation with the precise finite semantics claimed.

Carneades contributes a complementary account of ordinary premises,
assumptions, exceptions and burdens of proof. This helps explain why an
unanswered challenge and a rebutted proposition should not have the same
status. The original paper does not claim that its proposed proof standards
adequately model legal proof standards \citep[secs.~3--5]{carneades2007}.
The later \href{https://github.com/carneades/carneades-4}{Carneades~4}
implementation has different semantics and numerical conventions. Its
default proof standard, weights and unsupported-statement labels must not
be transferred into Hong Kong law or read as factual negation. We should
borrow the premise distinctions immediately and assess an engine choice
separately.

\subsection{Specify the reading before compiling it}

Arguments require propositions with stable meanings. LegalRuleML provides
ways to associate alternative interpretations with their sources,
authorities and contexts, including temporal and jurisdictional scope.
Its treatment of competing readings of an income-year provision shows why
two formalizations may need to coexist. It also distinguishes prescriptive
rules, which impose duties, from constitutive rules, which establish legal
classifications \citep[secs.~3--4]{legalinterpretations2014}.

For the circular, the representation should therefore state who must do what,
for which activity and product, during which interval, under which condition
and subject to which exception. It must preserve whether a condition governs
confirmation, demonstration, third-party verification or a legal opinion.
These are local design requirements motivated by the disputed source, not
a claim that adopting the LegalRuleML schema determines its interpretation.

Logical English supplies another useful component: controlled predicate
templates and explicit conventions for binding variables. These can make a
proposed specification readable while retaining executable structure.
The language reduces selected ambiguities by restricting the syntax accepted
as a specification \citep[secs.~1--2]{logicalenglish2023}. It does not
establish that the chosen specification faithfully translates unrestricted
English. Its basic untyped syntax and treatment of negation also require care
when facts are incomplete. A noun used as a variable name is not itself a
type constraint, and failure to establish supervision must not automatically
become evidence that supervision did not occur.

Figure~\ref{fig:demonstration-readings} isolates one consequence of the
request-condition dispute. Both branches are proposed readings; the other
relevant facts are held fixed for the illustration. The question is whether
the demonstration duty is triggered, not whether the provider complies
overall.

\begin{figure}[htbp]
\centering
\begin{tikzpicture}[
  box/.style={draw=navy,rounded corners,fill=pale,align=center,
    text width=6.1cm,inner sep=8pt,font=\small},
  link/.style={-{Latex[length=2mm]},draw=navy,thick}]
\node[box,text width=11.8cm] (case) at (0,0)
  {Declared situation: no SFC request has been received.\\
   Paragraph~14 and footnote~5 remain in the source.};
\node[box] (request) at (-3.65,-2.2)
  {Reading A: a request is required.\\[3pt]
   The request trigger is absent.};
\node[box] (standing) at (3.65,-2.2)
  {Reading B: the footnote supports a standing demonstration duty.\\[3pt]
   No request is needed to trigger it.};
\draw[link] (case.south) -- (request.north);
\draw[link] (case.south) -- (standing.north);
\end{tikzpicture}
\caption{The same declared fact has different consequences under rival
readings. The diagram isolates their difference; it does not decide which
reading the source warrants.}
\label{fig:demonstration-readings}
\end{figure}

A demonstration of this difference is useful even while its legal resolution
remains open. It tells the implementation exactly which choice must not be
hidden in a default. Producing a controlled-English sentence and executable
expression from the same record can then prevent them from drifting apart,
while a separate source challenge continues to examine the record itself.

\subsection{Evidence has to establish the right predicate}

A precise rule is still conditional on its facts. Holzenberger and Van Durme
connect legal information extraction with symbolic statutory reasoning by
representing events and their arguments through source-linked spans.
The representation ties a value to the passage from which it was obtained,
allowing downstream reasoning to expose the consequences of extraction errors.
The study also shows why a missing event argument can matter to the resulting
legal calculation \citep[secs.~3--5 and apps.~B, C and E]{connecting2023}.

That method gives us a useful evidence structure, with limits that should
shape its reuse. The study uses a tax-specific ontology and curated formal
statutes. Its date-completion conventions and closed-world treatment of
exceptions cannot silently become rules for incomplete regulatory evidence.
The associated implementation was inspected as source, not executed here;
its repository reuse licence was not established. The immediate proposal is
to implement the published span-linked idea within our own typed evidence
records.

For a judgment, the record should also say whether a passage reports an
allegation, a party's submission, a factual finding or judicial reasoning.
For an operational record, it should identify the provider, arrangement
component, event and time. A correctly quoted supervision log concerning
Provider A does not establish supervision of Provider B. A record entered
after the assessment cutoff was not available to the earlier assessment,
even if it describes an earlier event.

Positive and negative support must be kept separately. There may be support
for a proposition, support against it, both, or neither. These are states of
the declared evidence, not guarantees that its contents are true. A complete
record of inspections might support an inference that no inspection occurred
during a specified interval; an incomplete folder cannot support the same
inference merely because no inspection is found. The completeness premise
belongs in the argument.

Legal classifications require a further warrant. A blockchain record may
establish recorded token possession under its stated identity assumptions.
The conclusion that this proves legal title additionally needs the relevant
legal and contractual rules. Likewise, a recorded supervisory act need not
establish adequate supervision. The product should expose this distinction
between an observation it can derive and an evaluative standard it has not
established.

\subsection{A precise account of disagreement is an executable result}

The first implementation should evaluate the consequences of each retained
reading under explicitly declared facts. Let \(\mathcal R\) be the nonempty
finite set of retained reading branches, including their semantic settings
and assumptions. Let \(x\) be the scoped factual record. For each branch
\(r\), let \(d_r(x)\) return its declared duty outcome, retaining unknown,
conflicting or unsupported outcomes where appropriate. The set
\begin{equation}
  \mathcal D_{\mathcal R}(x)
    = \{\,d_r(x) : r\in\mathcal R\,\}
  \label{eq:retained-reading-outcomes}
\end{equation}
is a local definition of the outcomes reported across those branches. The
full report also retains each branch's outcome and premises; the set removes
duplicate values only for this comparison.

If every branch is evaluated and this set consists of exactly one settled
duty outcome, the evaluated branches agree on that outcome. A settled result
alongside an unknown result is not agreement. If it contains two distinct
settled outcomes, they disagree. If a branch remains unsupported, its status
must remain visible rather than being dropped to obtain agreement.
Equation~\eqref{eq:retained-reading-outcomes} therefore organizes a conditional
report; it does not show that \(\mathcal R\) contains every legally possible
reading or that a common interpretation is correct. A resource limit or an
unvisited alternative prevents a claim of completed evaluation.

For the no-request illustration, the two branches report different trigger
statuses. The product can establish that difference relative to the specified
rules while leaving the source question unresolved. With incomplete evidence
about whether a request was made, it must preserve a further factual
uncertainty. Keeping the two uncertainties separate tells the next
investigation whether it needs a record of events, an applicable authority,
or a repair to the representation.

Existing RuleIR, Catala and formal checking mechanisms can help verify these
declared consequences. They cannot supply the missing relationship between
the original English and the proposed formal specification. That relationship
must remain explicit even when multiple executable targets agree.

\subsection{The first repair program}

The proposed program begins with the two clause groups, preserving the
current source editions and every original obligation. It first corrects
demonstrably stale operative wording. It then builds the typed representation,
constructs the rival clause readings and evaluates their arguments under
explicit semantics. Precedent analysis adds opinion-aware propositions and
factor comparisons; evidence derivation adds actor, activity and time to the
facts the program actually uses.

The next stage joins explanation and execution inside the complete
investigator. Unsupported readings remain in its accounting. Independent
finite checks and narrowly stated formal propositions test preservation of
branches, evidence states and dependencies. Only then does the program extend
coverage and freeze a method for observations on genuinely later sources.
Each phase rewrites the next phase's plan in light of its actual results.
A counterexample normally identifies the next repair; corrupted evidence or
an invalid checker stops the affected execution.

This is proposed product work. No phase of that implementation program was
executed for this literature and document review. The practical dependency
choices are also bounded: PyArg, Logical English, Carneades and eyecite have
identified open-source licences; SARA-IE's code reuse licence and CLERC's
repository-wide coverage remain unverified. The retained source snapshots and
licence findings accompany the detailed program.\footnote{See
\srcpath{docs/plans/legal-interpretation-master-program.md} and the individual
plans under \srcpath{docs/plans/legal-interpretation-program/}. The source review
and pinned versions are in
\srcpath{docs/research/legal-interpretation-reuse-2026-10-05.md}.}

The resulting product should explain why a conclusion follows under a stated
reading and what would change it. Under the controlling research requirement,
human answer keys, reviewer approval and agreement among language models do
not qualify legal correctness. Published legal texts remain premises, and an
independently checked formal proposition establishes only its stated domain.
Where the source, applicability or evidence remains unsettled, the appropriate
result is a specific qualification accompanied by the consequences that can
already be established.
